\subsection{The weighted lattice sum for a fixed flag normal form}
\label{sec:weighted-flag-lattice-sum}

The integral transport theorem now supplies the lattice
representatives used in the arithmetic count.  We retain
the original common unipotent group and its stabilizer
weight.  The main cancellation is between the index of
the common isotropic lattice and the discriminant of the
paired flags.

Continue with the data of
Section~\ref{sec:flag-lattice-transport}, and put
\[
        q=2r+1,\qquad c_E=1+\log_2(2f),\qquad
        Q=D_b\Delta N_T.
\]
For an admissible lattice $\Lambda$, let
\[
 \mathcal J=(\Lambda+X)\cap(W+Ez),\qquad
 \mathcal J'=(\Lambda+X')\cap(W+Ez),\qquad
 L=\mathcal J\cap W,\quad L'=\mathcal J'\cap W.
 \tag{WF1}
\]
The actual congruence graph in $W\times W$ consists of
the pairs $(w',w)\in L'\times L$ admitting lifts
$\lambda'-w'\in X'$ and $\lambda-w\in X$, with
$\lambda',\lambda\in\Lambda$ and
\[
                 \lambda-\lambda'
                         \in g\Lambda+(\Lambda\cap X).
 \tag{WF2}
\]
Write $K$ for its diagonal fixed lattice.  The quotient
of the original theta lattice by $I$ is
\[
          M_0=\operatorname{pr}_W(\Lambda\cap(W+I)).
\]
Then $M_0\subset K\subset L$, and the original indices
are
\[
                     x=[L:K],\qquad i=[K:M_0].
 \tag{WF3}
\]
The common-isotropic index $b$ is as in (US5).

\begin{lemma}[The full integral index constraint]
\label{lem:flag-fixed-index-constraint}
For every original admissible lattice,
\[
                            xi^2b^2\leq\Delta.
 \tag{WF4}
\]
Moreover $\mathcal J$ and $\mathcal J'$ are unimodular;
their intersections with $Ez$ are respectively
$\mathcal O z$ and $\mathcal O cz$.
The definite lattice $L$ has cyclic discriminant module
$\mathcal O/T\mathcal O$, of order $N_T^2$.
\end{lemma}

\begin{proof}
We use the elementary integral isotropic-quotient fact:
if $A$ is a primitive isotropic submodule of a unimodular
Hermitian lattice, then $A^\perp/A$ is unimodular.
Here is a proof valid for nonfree modules.  Projectivity
and unimodularity give a lift of the Hermitian dual of
$A$ whose pairing with $A$ is perfect.  Subtract from
this lift the form $w_0$ times its own Gram form, using
the perfect pairing with $A$ to regard that form as a
map into $A$.  Since $w_0+\bar w_0=1$, the corrected
lift is isotropic.  It and $A$ form a unimodular
hyperbolic summand.  Its orthogonal complement is the
stated quotient and is unimodular.  This construction
uses the integral trace-one element already proved to
exist, so it also applies above $2$.

Apply this fact first to $\Lambda\cap I$.  In the
resulting unimodular lattice $\Lambda_I$, let
$\overline U,\overline U'$ be the images of
$\Lambda\cap X,\Lambda\cap X'$.  They are primitive
isotropic lattices of rank $j$ in the paired space
$H_j=X_j+X'_j$.  Let
\[
 B=\overline U+\overline U',\qquad
 P=\Lambda_I\cap H_j,\qquad
 C=\Lambda_I\cap(W+Ez),\qquad a=[P:B].
\]
The images of $\Lambda\cap X_j$ and
$\Lambda\cap X'_j$ have indices $b$ and $b'$ in
$\overline U$ and $\overline U'$.  These two indices
are equal.  Indeed projection to $I$ along the first
$j$ vectors has leading ideals $A_j,\ldots,A_{r-1}$
for the first flag and $cA_j,\ldots,cA_{r-1}$ for
the second.  Their covolumes coincide because
$|N_{E/\mathbf Q}c|=1$, and both contain the same
lattice $\Lambda\cap I$.

Let $\delta$ denote normalized Hermitian covolume, as
defined in the proof of (EI3).  The cross Gram determinant
of the two first-$j$ lattice flags gives
\[
        \delta(B)=\frac{\Delta}{b^2},\qquad
        \delta(P)=\frac{\Delta}{ab^2}.
 \tag{WF5}
\]
The sublattices $P,C$ are primitive orthogonal complements
in a unimodular lattice, so their discriminant modules
are paired by an anti-isometry.  In particular
$\delta(C)=\delta(P)$.

The isotropic-quotient fact applied to $\overline U$
and $\overline U'$ identifies the two quotient lattices
with $\mathcal J$ and $\mathcal J'$ in (WF1).  Both
are unimodular and contain $C$.  Therefore
\[
                         n=[\mathcal J:C]
                                =\frac{\Delta}{ab^2}.
 \tag{WF6}
\]
There is also the bound
\[
                         [\mathcal J\cap\mathcal J':C]\leq a.
 \tag{WF7}
\]
To verify it, use the discriminant-module gluing of
$P\perp C$ to $\Lambda_I$.  The subgroups corresponding
to $\mathcal J/C$ and $\mathcal J'/C$ consist of the
classes in $P^\#/P$ represented in the two isotropic
spaces.  A common class has representatives $u,u'$ in
those spaces, with $u-u'\in P$.  Sending the class to
$u-u'$ modulo $B$ is well-defined: changing either
representative changes it by an element of
$P\cap X_j=\overline U$ or
$P\cap X'_j=\overline U'$, respectively.
It is injective into $P/B$.  If the image
is zero, uniqueness of the two isotropic components
puts $u\in\overline U$ and $u'\in\overline U'$, so
the class was zero.  This proves (WF7).

Now $M_0=C\cap W$.  Thus $L/M_0$ injects into
$\mathcal J/C$, and $K/M_0$ injects into
$(\mathcal J\cap\mathcal J')/C$.
Equations (WF3), (WF6) and (WF7) give $xi\leq n$ and
$i\leq a$.  Multiplication proves (WF4).

Finally the leading conditions on the last vectors give
$\mathcal J\cap Ez=\mathcal O z$ and
$\mathcal J'\cap Ez=\mathcal O cz$.
The first is a primitive rank-one lattice of norm $-T$
inside a unimodular lattice.  Discriminant gluing with
its primitive orthogonal complement $L$ identifies
$L^\#/L$ with the opposite discriminant module of
$\mathcal O z$, namely $\mathcal O/T\mathcal O$.
This proves the last assertion.
\end{proof}

For a lattice $B$, write $\Gamma_B=R\cap\operatorname{Stab}(B)$.
The weight in the source is
\[
           \omega(\Lambda)=
           \operatorname{ivol}(\Gamma_{\Lambda_{\mathrm{ref}}},
                               \Gamma_\Lambda).
 \tag{WF8}
\]
Fix the two actual lattices $\mathcal J,\mathcal J'$
and the graph (WF2).  In particular $L,L',K$ and $x$
are fixed.  The weight sum below runs over common
unipotent $R$-orbits of integral unimodular lattices
with precisely those data and the prescribed flag and
congruence conditions.

\begin{theorem}[The original fixed-normal-form weighted sum]
\label{thm:fixed-normal-form-weighted-sum}
For specified positive integers $i,b$, the weight sum
over lattices having those indices is at most
\[
                  (ib)^k i^{c_E}Q^{10000(r+1)^3}.
 \tag{WF9}
\]
Without fixing $i,b$, it is at most
\[
              (\Delta/x)^{k/2}\Delta^{c_E+2}
                                     Q^{10000(r+1)^3}.
 \tag{WF10}
\]
If only $\mathcal J,L'$ and the congruence graph are
fixed, multiply the right side of (WF10) by $N_T^2$.
These are the original quantitative bounds, with their
original stabilizer weights.
\end{theorem}

\begin{proof}
The fixed data and indices are unchanged by $R$.
For instance its action on $(X+X')^\perp/I$ is the
identity, while its action on each flag is the identity
on the successive quotients.  The descriptions (WF1)
and (WF2) are consequently preserved.  The index $b$
is unchanged because $X,X_j,I$ are preserved.

The weight is also invariant under this action.  Here
one can avoid a Haar normalization entirely.  Stabilizers
of any two full lattices in a rational linear group are
commensurable: a sufficiently deep integer congruence
subgroup for either lattice preserves the other.
For such a stabilizer $\Gamma$, the relative-index
modulus
\[
       u\longmapsto
          \operatorname{ivol}(\Gamma,u\Gamma u^{-1})
\]
is a homomorphism to $\mathbf Q_{>0}$, by the finite-index
multiplication rule.  Every element of the rational
unipotent group $R$ has an $n$th root in $R$, for every
$n\geq1$, given by the finite binomial polynomial in
$u-1$.  That polynomial remains unitary and preserves
both flag filtrations: its adjoint is its inverse,
and its differences have the same strictly triangular
containments.  The displayed positive rational number
is therefore an $n$th power in $\mathbf Q$ for every
$n$.  Its prime valuations are all zero, so it is $1$.
Apply this to $\Gamma=\Gamma_\Lambda$ and use
$\Gamma_{u\Lambda}=u\Gamma_\Lambda u^{-1}$.
The relative-index multiplication rule in (WF8)
now proves weight invariance.

Apply Theorem~\ref{thm:flag-lattice-near-representative}
with $e=12r+50$.  Every orbit has a split representative
whose lattice $P=\Lambda\cap\Pi$ satisfies
$\mathfrak c^eP_0\subset P\subset\mathfrak c^{-e}P_0$.
For such a representative, $M=\Lambda\cap W=M_0$.
There are at most $i^{k+c_E}$ choices for $M\subset K$
of index $i$, by the rank-$k$ sublattice count proved
with (MM13).  The finite module
\[
                     \mathfrak c^{-e}P_0/\mathfrak c^eP_0
\]
has $(N\mathfrak c)^{2eq}$ elements and at most $q$
generators.  Every submodule also has at most $q$
generators, by localization and the Chinese remainder
theorem.  Thus there are at most
$(N\mathfrak c)^{2eq^2}$ choices for $P$.

For fixed $M,P$, unimodularity makes $\Lambda/(M\perp P)$
the graph of an anti-isometry between their full
discriminant modules.  Conversely such a graph determines
$\Lambda$.  The module $P^\#/P$ has at most $q$
generators.  Since (WF4) gives $xi\leq\Delta$,
\[
                   \Delta(M)=N_T^2(xi)^2
                              \leq(\Delta N_T)^2.
\]
There are consequently at most
$(\Delta N_T)^{2q}$ gluings.  This is only an upper
bound, so retaining all anti-isometries can enlarge
the count without discarding any original congruence.

The stabilizer estimate (FT15) bounds each representative's
weight by
$b^k\Delta^r(N\mathfrak c)^{2(3e+11)q^2}$.
Multiplying the three counts gives
\[
 \sum\omega(\Lambda)
 \leq (ib)^k i^{c_E}\Delta^r
       (\Delta N_T)^{2q}
       (N\mathfrak c)^{(8e+22)q^2}.
 \tag{WF11}
\]
Here $N\mathfrak c=(D_bN_T)^2\leq Q^2$ and every
factor $D_b,\Delta,N_T$ is at least $1$.  For
$q=2r+1$ and $e=12r+50$,
\[
 r+2q+2(8e+22)q^2
       =768r^3+4144r^2+3573r+846
       \leq10000(r+1)^3.
\]
This proves (WF9) and in particular finiteness of the
set of relevant orbits.

For all allowed pairs $(i,b)$, (WF4) gives
$(ib)^k\leq(\Delta/x)^{k/2}$ and $i^{c_E}\leq\Delta^{c_E}$.
There are at most $\Delta^2$ such pairs.  Summing (WF9)
proves (WF10).

To remove the specification of $\mathcal J'$, note that
it is a unimodular gluing of the fixed $L'$ and the
fixed rank-one lattice $\mathcal O cz$.  The latter's
discriminant module is cyclic of order $N_T^2$.
An anti-isometry from it is determined by the image of
one generator, so there are at most $N_T^2$ possible
gluings.  This proves the final assertion.
\end{proof}

\paragraph{The finite normal-form parameters.}
For a prescribed leading-minor label, the ordered signed
block antidiagonal matrix in (FT1) has at most $8^j$
possibilities.  There are at most $2^j$ sign choices
and $2^{j-1}$ consecutive-block decompositions.  For a
block of length $d$, its scalar satisfies an equation
of degree $d$, determined by the adjacent leading
determinants; it therefore has at most $d$ choices.
The product of the block lengths is at most $2^j$.
This proves the stated bound, including the empty
matrix when $j=0$.

The scalar $c$ with $c\bar c=1$ also has a uniform
finite bound.  Define the fractional ideal
\[
                  C_z=\{t\in E:tz\in\Lambda+X+X'\}.
\]
Both $\mathcal O$ and $c\mathcal O$ lie in $C_z$.
The paired Gram determinant shows
$\mathfrak a_jC_z\subset\mathcal O$ or the same
inclusion with $\overline{\mathfrak a_j}$.
Indeed, in a lattice lift $tz+x+x'$, integrality
against $\Lambda\cap X_j$ controls the $X'_j$
component; the adjugate puts its product with the
terminal determinant ideal into $\Lambda\cap X'_j$.
Subtracting it leaves a first-flag lift, whose last
leading coefficient is integral.
Consequently
\[
        (c)=\mathfrak b\mathfrak a^{-1},\qquad
        \mathfrak a=C_z^{-1},\quad
        \mathfrak b=cC_z^{-1},\qquad
        N\mathfrak a=N\mathfrak b\leq\Delta.
 \tag{WF12}
\]
Both ideals are integral.  For fixed ideals, two such
scalars differ by a unit all of whose conjugates have
absolute value $1$, hence by a root of unity.
Indeed the characteristic polynomials over $\mathbf Q$
of all positive powers of such a unit have integral
coefficients bounded by the binomial coefficients in
degree $2f$.  Only finitely many polynomials, and hence
finitely many complex roots, occur; two powers coincide.
The
elementary inequality $\varphi(n)^2\geq n/2$ shows
that $E$ has at most $8f^2$ roots of unity: prove the
inequality prime-power by prime-power, with the only
factor below $1$ being $1/2$ for $2^1$.
The ideal-count bound already used above gives at most
$\Delta^{1+\log_2(2f)}$ ideals of norm at most $\Delta$.
Thus the number of scalars in (WF12) is at most
\[
                     8f^2\Delta^{2(2+c_E)}.
 \tag{WF13}
\]
The exponent stated here is the original, looser one.
Combining (WF10), the $N_T^2$ extension count, (WF13)
and the signed block count gives the original finite-data
bound
\[
 (\Delta/x)^{k/2}\,
 8^{r+1}f^2\Delta^{3c_E+6}N_T^2
                   (D_b\Delta N_T)^{10000(r+1)^3}.
 \tag{WF14}
\]
This expression is precisely \texttt{datBd} in the source's
\texttt{AdaptSelf}, lines 42--44: its $d_b,d,x$ are
our $D_b,\Delta,x$, its $\mathrm{NT}$ is $N_T$, and
its $k$ is $\dim_EW$~\cite{weighted-flag-lattice-source}.
To use (WF14) as the full per-row term bound, one must
still verify the original normal-form assignment,
the equivalence of flag classes under its transport,
and the reciprocal row-automorphism multiplicity.
The fixed-normal-form arithmetic and its covolume
weight are fully supplied by (WF9)--(WF13); those
remaining term interfaces are separate obligations.
